%% abtex2-modelo-trabalho-academico.tex, v<VERSION> laurocesar
%% Copyright 2012-<COPYRIGHT_YEAR> by abnTeX2 group at http://www.abntex.net.br/ 
%%
%% This work may be distributed and/or modified under the
%% conditions of the LaTeX Project Public License, either version 1.3
%% of this license or (at your option) any later version.
%% The latest version of this license is in
%%   http://www.latex-project.org/lppl.txt
%% and version 1.3 or later is part of all distributions of LaTeX
%% version 2005/12/01 or later.
%%
%% This work has the LPPL maintenance status `maintained'.
%% 
%% The Current Maintainer of this work is the abnTeX2 team, led
%% by Lauro César Araujo. Further information are available on 
%% http://www.abntex.net.br/
%%
%% This work consists of the files abntex2-modelo-trabalho-academico.tex,
%% abntex2-modelo-include-comandos and abntex2-modelo-references.bib
%%

% ------------------------------------------------------------------------
% ------------------------------------------------------------------------
% abnTeX2: Modelo de Trabalho Academico (tese de doutorado, dissertacao de
% mestrado e trabalhos monograficos em geral) em conformidade com 
% ABNT NBR 14724:2011: Informacao e documentacao - Trabalhos academicos -
% Apresentacao
% ------------------------------------------------------------------------
% ------------------------------------------------------------------------

\documentclass[
	% -- opções da classe memoir --
	12pt,				% tamanho da fonte
	openright,			% capítulos começam em pág ímpar (insere página vazia caso preciso)
	twoside,			% para impressão em recto e verso. Oposto a oneside
	a4paper,			% tamanho do papel. 
	% -- opções da classe abntex2 --
	%chapter=TITLE,		% títulos de capítulos convertidos em letras maiúsculas
	%section=TITLE,		% títulos de seções convertidos em letras maiúsculas
	%subsection=TITLE,	% títulos de subseções convertidos em letras maiúsculas
	%subsubsection=TITLE,% títulos de subsubseções convertidos em letras maiúsculas
	% -- opções do pacote babel --
	english,			% idioma adicional para hifenização
	french,				% idioma adicional para hifenização
	spanish,			% idioma adicional para hifenização
	brazil				% o último idioma é o principal do documento
	]{abntex2}

% ---
% Pacotes básicos 
% ---
\usepackage{lmodern}			% Usa a fonte Latin Modern			
\usepackage[T1]{fontenc}		% Selecao de codigos de fonte.
\usepackage[utf8]{inputenc}		% Codificacao do documento (conversão automática dos acentos)
\usepackage{lastpage}			% Usado pela Ficha catalográfica
\usepackage{indentfirst}		% Indenta o primeiro parágrafo de cada seção.
\usepackage{color}				% Controle das cores
\usepackage{graphicx}			% Inclusão de gráficos
\usepackage{microtype} 			% para melhorias de justificação
% ---
% Adicionado por Madson Dias
% ---

\usepackage{utils/ifce-modelo} 					% Pacote com informações da customização
\graphicspath{ {figuras/graficos/} {figuras/imagens/}}  % Inclusão dos paths para imagens
\usepackage{xstring} 									% Criar comandos com IF
\usepackage{caption}
\usepackage{subcaption}
% ---
% Pacotes adicionais, usados apenas no âmbito do Modelo Canônico do abnteX2
% ---
\usepackage{lipsum}				% para geração de dummy text
% ---

% ---
% Pacotes de citações
% ---
\usepackage[brazilian,hyperpageref]{backref}	 % Paginas com as citações na bibl
\usepackage[alf]{abntex2cite}	% Citações padrão ABNT

% --- 
% CONFIGURAÇÕES DE PACOTES
% --- 

% ---
% Configurações do pacote backref
% Usado sem a opção hyperpageref de backref
\renewcommand{\backrefpagesname}{Citado na(s) página(s):~}
% Texto padrão antes do número das páginas
\renewcommand{\backref}{}
% Define os textos da citação
\renewcommand*{\backrefalt}[4]{
	\ifcase #1 %
		Nenhuma citação no texto.%
	\or
		Citado na página #2.%
	\else
		Citado #1 vezes nas páginas #2.%
	\fi}%
% ---

% ---
% Informações de dados para CAPA e FOLHA DE ROSTO
% ---
\titulo{Uso de metaheuristicas no treinamento de Máquinas de Vetores Suporte}
\autor{Madson Luiz Dantas Dias}
\local{Fortaleza, CE}
\data{Março, 2015}
\orientador{Ajalmar Rego Rocha Neto}
\coorientador{Pedro Pedrosa Rebouças Filho}
\instituicao{%
  Instituto Federal do Ceará (IFCE)
  \par
  Programa de Pós-Graduação em Ciência da Computação (PPGCC)}
\tipotrabalho{Dissertação de Mestrado}
% O preambulo deve conter o tipo do trabalho, o objetivo, 
% o nome da instituição e a área de concentração 
\preambulo{Tese submetida ao Programa de Pós-Graduação em Ciência da Computação, do Instituto Federal do Ceará, como requisito parcial para obtenção do grau de Mestre em Ciência da Computação

Área de Concentração: Inteligência Artificial.

Orientador: Ajalmar Rego Rocha Neto.}

% ---


% ---
% Configurações de aparência do PDF final

% alterando o aspecto da cor azul
\definecolor{blue}{RGB}{41,5,195}

% informações do PDF
\makeatletter
\hypersetup{
     	%pagebackref=true,
		pdftitle={\@title}, 
		pdfauthor={\@author},
    	pdfsubject={\imprimirpreambulo},
	    pdfcreator={LaTeX with abnTeX2},
		pdfkeywords={abnt}{latex}{abntex}{abntex2}{trabalho acadêmico}, 
		colorlinks=true,       		% false: boxed links; true: colored links
    	linkcolor=blue,          	% color of internal links
    	citecolor=blue,        		% color of links to bibliography
    	filecolor=magenta,      		% color of file links
		urlcolor=blue,
		bookmarksdepth=4
}
\makeatother
% --- 

% --- 
% Espaçamentos entre linhas e parágrafos 
% --- 

% O tamanho do parágrafo é dado por:
\setlength{\parindent}{1.3cm}

% Controle do espaçamento entre um parágrafo e outro:
\setlength{\parskip}{0.2cm}  % tente também \onelineskip

% ---
% compila o indice
% ---
\makeindex
% ---

% ----
% Início do documento
% ----
\begin{document}

% Seleciona o idioma do documento (conforme pacotes do babel)
%\selectlanguage{english}
\selectlanguage{brazil}

% Retira espaço extra obsoleto entre as frases.
\frenchspacing 

% ----------------------------------------------------------
% ELEMENTOS PRÉ-TEXTUAIS
% ----------------------------------------------------------
% \pretextual
\imprimircapa % Capa *
\imprimirfolhaderosto % Folha de rosto *

\begin{folhadeaprovacao}

    \begin{center}
    \includegraphics[width=2.5cm]{brasao_republica.jpg} \\
    \vspace{.5cm}
    {\ABNTEXchapterfont\large\imprimirinstituicao}
    \vspace{1cm}

    {\ABNTEXchapterfont\large\imprimirautor}

    % \vspace*{\fill}\vspace*{\fill}
    % \begin{center}
    %   \ABNTEXchapterfont\bfseries\Large\imprimirtitulo
    % \end{center}
    %\vspace*{\fill}
    
    %\hspace{.45\textwidth}
    % \begin{minipage}{.5\textwidth}
    %     \imprimirpreambulo
    %     {\large\imprimirdata}
    % \end{minipage}%
    %\vspace*{\fill}
   \end{center}
   Esta dissertação foi julgada adequada para a obtenção do título de Mestre em Ciência da
Computação, sendo aprovada pela Coordenação do Programa de Pós-Graduação em Ciência
da Computação do Instituto Federal de Educação, Ciência e Tecnologia do Ceará e pela banca
examinadora:

   \assinatura{\textbf{\imprimirorientadorRotulo~\imprimirorientador} \\ Instituto Federal de Educação, Ciência e Tecnologia do Ceará -- IFCE}

   \if \imprimircoorientador
    
   \else
    \assinatura{\textbf{\imprimircoorientadorRotulo~\imprimircoorientador}  \\ Instituto Federal de Educação, Ciência e Tecnologia do Ceará -- IFCE}
   \fi

   \if \nomeprofessorA \instituicaoprofessorA
    
   \else
    \assinatura{\textbf{\imprimirnomeprofessorA}  \\ \imprimirinstituicaoprofessorA}
   \fi

   \if \nomeprofessorB
    
   \else
    \assinatura{\textbf{\imprimirnomeprofessorB}  \\ \imprimirinstituicaoprofessorB}
   \fi
   
   \if \nomeprofessorC
    
   \else
    \assinatura{\textbf{\imprimirnomeprofessorC}  \\ \imprimirinstituicaoprofessorC}
   \fi
   
   
   %\assinatura{\textbf{Professor}  \\ Instituto Federal do Ceará (IFCE)}
      
   \begin{center}
    \vspace*{\fill}
    {\large\imprimirlocal}
    \par
    {\large\imprimirdata}
    \vspace*{1cm}
  \end{center}
  
\end{folhadeaprovacao}


%\includepdf{pretextual/folhadeaprovacao.pdf} 	% Folha de aprovação (Depois da apresentação)
\begin{dedicatoria}
   \vspace*{\fill}
   \centering
   \noindent
   \textit{ Este trabalho é dedicado às crianças adultas que,\\
   quando pequenas, sonharam em se tornar cientistas.} \vspace*{\fill}
\end{dedicatoria} 				% Dedicatória
\begin{agradecimentos}
Os agradecimentos principais são direcionados à Gerald Weber, Miguel Frasson,
Leslie H. Watter, Bruno Parente Lima, Flávio de Vasconcellos Corrêa, Otavio Real
Salvador, Renato Machnievscz\footnote{Os nomes dos integrantes do primeiro
projeto abn\TeX\ foram extraídos de
\url{http://codigolivre.org.br/projects/abntex/}} e todos aqueles que
contribuíram para que a produção de trabalhos acadêmicos conforme
as normas ABNT com \LaTeX\ fosse possível.

Agradecimentos especiais são direcionados ao Centro de Pesquisa em Arquitetura
da Informação\footnote{\url{http://www.cpai.unb.br/}} da Universidade de
Brasília (CPAI), ao grupo de usuários
\emph{latex-br}\footnote{\url{http://groups.google.com/group/latex-br}} e aos
novos voluntários do grupo
\emph{\abnTeX}\footnote{\url{http://groups.google.com/group/abntex2} e
\url{http://www.abntex.net.br/}}~que contribuíram e que ainda
contribuirão para a evolução do \abnTeX.

\end{agradecimentos} 			% Agradecimento
\begin{epigrafe}
    \vspace*{\fill}
	\begin{flushright}
		\textit{%
		``<Citação Célebre>''\\
		(<Autor da citação>)}
	\end{flushright}
\end{epigrafe} 				% Epígrafe
\setlength{\absparsep}{18pt} % ajusta o espaçamento dos parágrafos do resumo
\begin{resumo}
 Segundo a \citeonline[3.1-3.2]{NBR6028:2003}, o resumo deve ressaltar o
 objetivo, o método, os resultados e as conclusões do documento. A ordem e a extensão
 destes itens dependem do tipo de resumo (informativo ou indicativo) e do
 tratamento que cada item recebe no documento original. O resumo deve ser
 precedido da referência do documento, com exceção do resumo inserido no
 próprio documento. (\ldots) As palavras-chave devem figurar logo abaixo do
 resumo, antecedidas da expressão Palavras-chave:, separadas entre si por
 ponto e finalizadas também por ponto.

 \textbf{Palavras-chave}: latex. abntex. editoração de texto.
\end{resumo} 					% Resumo em português 
\begin{resumo}[Abstract]
 \begin{otherlanguage*}{english}
   This is the english abstract.

   \vspace{\onelineskip}
 
   \noindent 
   \textbf{Keywords}: latex. abntex. text editoration.
 \end{otherlanguage*}
\end{resumo} 				% Resumo em inglês 
% inserir lista de ilustrações
\pdfbookmark[0]{\listfigurename}{lof}
\listoffigures*
\cleardoublepage
% inserir lista de tabelas
\pdfbookmark[0]{\listtablename}{lot}
\listoftables*
\cleardoublepage
% ---
\begin{siglas}
  \item[ABNT] Associação Brasileira de Normas Técnicas
  \item[abnTeX] ABsurdas Normas para TeX
\end{siglas}  			% Lista de abreviaturas e siglas
\begin{simbolos}
\item[$ \Gamma $] Letra grega Gama
\item[$ \Lambda $] Lambda
\item[$ \zeta $] Letra grega minúscula zeta
\item[$ \in $] Pertence
\end{simbolos} 				% Lista de símbolos
% inserir o sumario
\pdfbookmark[0]{\contentsname}{toc}
\tableofcontents*
\cleardoublepage
% ---



% ----------------------------------------------------------
% ELEMENTOS TEXTUAIS
% ----------------------------------------------------------
\textual

\chapter{Introdução}\label{cap:introducao}
\lipsum[2-3]
\lipsum[1-4]
\lipsum[3-6]

\section{Motivação para a Dissertação}\label{sec:motivacao}
\lipsum[5-6]
\lipsum[1-2]
\lipsum[3-4]

\section{Descrição do problema}\label{sec:descricao-problema}
\lipsum[3-4]
\lipsum[1-2]
\lipsum[3-4]
\lipsum[1-2]

\section{Objetivos Geral e Específicos}\label{sec:objetivos}
\lipsum[1]

Em seu sentido mais particular, os seguintes objetivos específicos são:

\begin{itemize}
	\item Excepteur sint occaecat cupidatat non proident, sunt in culpa qui officia deserunt mollit anim id est laborum;
	\item Lorem ipsum dolor sit amet, consectetur adipisicing elit, sed do eiusmod;
	\item Tempor incididunt ut labore et dolore magna aliqua. Ut enim ad minim veniam;
	\item Quis nostrud exercitation ullamco laboris nisi ut aliquip ex ea commodo consequat. Duis aute irure dolor in reprehenderit in voluptate velit esse;
	\item Cillum dolore eu fugiat nulla pariatur. Excepteur sint occaecat cupidatat non.
\end{itemize}

\section{Produção científica}\label{sec:producao}
Durante este projeto de mestrado, os seguintes trabalhos científicos foram aceitos e publicados, a saber:

\begin{itemize}
	\item \textbf{Einstein, A.}, 1905. \textbf{The photoelectric effect}. Ann. Phys, 17(132), p.4;
	\item \textbf{Einstein, A.}, 1904. \textbf{Zur allgemeinen molekularen Theorie der Wärme}. Annalen der Physik, 319(7), pp.354-362.
\end{itemize}

\section{Estrutura da Dissertação}\label{sec:estrutura}
\lipsum[4-5]
\lipsum[1-2]
\lipsum[1-4]
\lipsum[3-2] 				% Capítulo de introdução
\chapter{Lorem ipsum dolor sit amet}\label{cap:exampleChapter}
% ---
\section{Aliquam vestibulum fringilla lorem}
% ---

\lipsum[1]

\figurasimples[talbot2012]{logo-ifce-sem-nome}{Exemplo de figura com fonte.}{4cm}

% ------
\lipsum[2-3]

\figurasimples{brasao_republica}{Exemplo de figura sem fonte.}{4cm}

\lipsum[2-3]

\tabela{tabela-exemplo}{Lista de bases de dados usados nesse trabalho}{1} % Exemplo de inclusão de tabela

\lipsum[2-3]

\lipsum[2-3]
% Exemplo de inclusão de gráficos

\figuradupla{pid}{Legenda da imagem 1}{rip}{Legenda da imagem 2}
% ------
\lipsum[2]

\begin{algorithm}
\caption{Algoritmo de exemplo}\label{bogosort}
  \begin{algorithmic}[1]
  \Procedure{Bogosort}{array}
    \While{\Not \Call{Está\_ordenado}{array}}
    \State array $\gets$ \Call{Permutação\_aleatória}{array} \Comment{Isso vai demorar muito :(}
    \EndWhile
  \EndProcedure
  \end{algorithmic}
\end{algorithm}

\lipsum[1] 			% Capitulo de exemplo

% ----------------------------------------------------------
% Finaliza a parte no bookmark do PDF
% para que se inicie o bookmark na raiz
% e adiciona espaço de parte no Sumário
% ----------------------------------------------------------
\phantompart

% ---
% Conclusão
% ---
\chapter{Conclusão}
% ---

\lipsum[31-33] 				% Capítulo de introdução

% ----------------------------------------------------------
% ELEMENTOS PÓS-TEXTUAIS
% ----------------------------------------------------------
\postextual
% ----------------------------------------------------------

% ----------------------------------------------------------
% Referências bibliográficas
% ----------------------------------------------------------
\bibliography{referencias}

% ----------------------------------------------------------
% Glossário
% ----------------------------------------------------------
%
% Consulte o manual da classe abntex2 para orientações sobre o glossário.
%
%\glossary

% ----------------------------------------------------------
% Apêndices
% ----------------------------------------------------------

% ---
% Inicia os apêndices
% ---
\begin{apendicesenv}

% Imprime uma página indicando o início dos apêndices
\partapendices
% Insere os apêndices
% ----------------------------------------------------------
\chapter{Nullam elementum urna vel imperdiet sodales elit ipsum pharetra ligula
ac pretium ante justo a nulla curabitur tristique arcu eu metus}
% ----------------------------------------------------------
\lipsum[55-57]
\chapter{Quisque libero justo}
% ----------------------------------------------------------

\lipsum[50]
%\include{apendices/apendice-c}
%\include{apendices/apendice-d}




\end{apendicesenv}
% ---


% ----------------------------------------------------------
% Anexos
% ----------------------------------------------------------

% ---
% Inicia os anexos
% ---
\begin{anexosenv}

% Imprime uma página indicando o início dos anexos
\partanexos
% Insere os anexos
\chapter{Morbi ultrices rutrum lorem.}
% ---
\lipsum[30]

\chapter{Cras non urna sed feugiat cum sociis natoque penatibus et magnis dis
parturient montes nascetur ridiculus mus}
% ---

\lipsum[31]
%\include{anexos/anexo-c}
%\include{anexos/anexo-d}

\end{anexosenv}

%---------------------------------------------------------------------
% INDICE REMISSIVO
%---------------------------------------------------------------------
\phantompart
\printindex
%---------------------------------------------------------------------

\end{document}
